\begin{frame}{확인 문제 (3)-(5)}
  \small
  \begin{itemize}
    \item \textbf{3.} ``GAE의 $\lambda=0.95$가 편향-분산을 조절해서
      학습이 안정됐다(Ch11.3)''로 끝난 정당화는 \textbf{2점에
      가깝다} — 개념은 인용했지만 \emph{관측 숫자가 없음}. 3점
      전환: ``$\lambda$를 0.8/1.0으로 바꿔 어드밴티지 \emph{추정치의
      분산}을 실측했습니다 — 실측 표가 보고에 있습니까?''(관찰+개념+
      숫자).
    \item \textbf{4.} ``왜 DQN이 아니라 PPO인가요?'' — 이대로면
      \textbf{2점}(``연속이라서'' 한 줄이면 반박이 어려움). 3점
      버전: ``Pendulum의 행동은 \emph{연속}인데 DQN(Ch09)은 \emph{이산}
      전제로 $Q(s,a)$를 출력 — 이산화의 \emph{정밀도 손실}을 확인했나요?
      아니면 연속 정책인 PPO(Ch10$\sim$11)가 자연스러운 이유가
      무엇인가요?''(체크리스트 2번에 Ch09 vs Ch10 개념 연결).
    \item \textbf{5.} ``PPO가 \emph{항상} 좋은가요? ``잘 학습했다''가
      무작위(약 $-1239$)보다 나은 걸 어떻게 아나요?'' — \textbf{정직한
      범위}: ``이 환경·이 설정($\gamma, \lambda$, 12만 스텝)에서 시드 3개
      한도에서 \emph{학습 도중} 이동평균이 무작위보다 높다''(단, 순수
      재평가 $-1413/-1491$로는 오히려 \emph{나쁨}; 외부 일반화는 이
      실험이 지지하지 않음). \textbf{말하면 안 될} 추론: ``PPO가 더 좋은
      알고리즘이다'' — 이 비교는 \emph{이 환경·이 스텝 수}의 답이지
      \textbf{알고리즘 우열}의 답이 아니다.
  \end{itemize}
\end{frame}
